\section{Results}

We evaluate detection quality on the held-out test split, generalization across recordings, per-leg attribution, probability calibration, feature importance, deployment robustness, and runtime cost. Unless stated otherwise, every number reported here comes from offline replay of recorded logs; we make no claim about live on-robot detection performance.

\subsection{Detection on the held-out test split}
Our first question is whether the model separates entangled windows from everything else better than a hand-tuned heuristic operating on the same signal. Table~\ref{tab:detection} compares the TCN against a rule-based baseline that thresholds thigh torque, evaluated under an identical protocol on the same five test recordings. On the test split the TCN reaches precision $0.799$, recall $0.818$, and $F1 = 0.809$, with $\text{PR-AUC} = 0.808$ and $\text{ROC-AUC} = 0.956$. The rule-based baseline attains comparable precision ($0.735$) but collapses on recall ($0.131$), for $F1 = 0.222$; a fixed torque threshold catches only the most violent snags and misses the many entanglements that manifest as sustained resistive effort rather than a single large spike. The gap in $F1$ ($0.809$ vs.\ $0.222$) and in exact-match multi-label accuracy ($0.908$ vs.\ $0.106$) is the central detection result.

The confusion matrix in Fig.~\ref{fig:confusion} makes the error profile concrete at the operating point (raw threshold $0.9999$): TP $2833$, FP $712$, FN $630$, TN $9126$. Errors are roughly balanced between false alarms and misses, and the large true-negative count reflects the dominance of ordinary walking and stationary windows in the data.

Because a single threshold hides how the score behaves across all thresholds, Fig.~\ref{fig:rocpr} plots the full ROC and precision-recall curves. The ROC-AUC of $0.956$ indicates strong ranking of positive over negative windows; the lower PR-AUC of $0.808$ is expected under class imbalance, where precision is sensitive to the many negatives.

\subsection{Generalization across recordings (LORO)}
Test-split numbers depend on one particular grouping of recordings, so we also run leave-one-recording-out (LORO) cross-validation, which retrains with each recording held out in turn. Across the 15 folds shown in Fig.~\ref{fig:loro}, $F1 = 0.807 \pm 0.142$. The mean tracks the test-split $F1$ closely, which suggests the split is representative rather than lucky. The spread is not uniform: the weakest folds are the front-right recording, the least represented leg, together with the deployment recordings that capture entanglement while the robot stands at rest, a regime that differs most from the walking-dominated remainder of the corpus.

\subsection{Per-leg attribution}
Detecting that \emph{some} leg is entangled is only half the task; the system must also name the leg. Fig.~\ref{fig:perleg} and Table~\ref{tab:perleg} report per-leg detection on the test split. Recall is high across all four legs (FR $1.000$, FL $1.000$, RR $0.959$, RL $0.916$), so genuine entanglements are rarely missed regardless of which leg is involved. The differences live in precision. Front-right shows a clear precision/recall trade-off: it recovers every front-right event (recall $1.000$) but at precision $0.467$, giving $F1 = 0.636$, the lowest of the four. This is consistent with the LORO finding, since FR has the fewest positive appearances and the model compensates by firing readily on FR. The other legs are stronger: FL $F1 = 0.845$ (P $0.731$, R $1.000$), RL $F1 = 0.823$ (P $0.748$, R $0.916$), and RR $F1 = 0.700$ (P $0.551$, R $0.959$).

\subsection{Calibration}
A detector that outputs probabilities should ideally have those probabilities mean what they say. Fig.~\ref{fig:reliability} reports the reliability diagram before and after temperature scaling. The raw network is badly over-confident: roughly $27\%$ of window probabilities are pinned near $1.0$, and this saturation forces the false-alarm-controlled threshold to clamp at $0.9999$. Temperature scaling with $T = 5.918$ de-saturates the probabilities and lets the operating threshold move to a more usable $0.963$. The effect on scoring metrics is mixed: the Brier score improves slightly, from $0.128$ to $0.118$, but the expected calibration error does \emph{not} improve, moving from $0.134$ to $0.156$. Temperature scaling here buys us a de-clamped, interpretable threshold and a marginally better Brier score, not better calibration in the ECE sense.

\subsection{Feature ablation}
To understand which input channels carry the signal, we retrain with feature groups removed and re-run LORO for each configuration. Table~\ref{tab:ablation} and Fig.~\ref{fig:ablation} report the change in F1 relative to the full 60-channel set. Raw kinematics are essential: using the engineered channels alone drops F1 by $0.13$, and removing all raw channels is by far the largest degradation. The engineered channels still earn their place, since raw-only sits $0.03$ below the full set. Foot contacts matter little (removing them costs $0.01$) and the IMU is effectively redundant for this task: dropping it leaves F1 unchanged to two decimals. In short, joint position, velocity, and torque do the work, the engineered snag-signature channels add a small margin, and the inertial channels contribute almost nothing.

\subsection{Deployment robustness}
Offline test metrics do not capture failure modes that appear on specific field recordings, so we retrained with targeted data augmentation and measured per-phase firing rates before (v1) and after (v2) on the relevant recordings; Fig.~\ref{fig:field} summarizes these. Four false-fire modes were driven to essentially zero: the Lock false-fire on \texttt{lock\_stop} ($0.8 \to 0.0$) and on \texttt{walk\_stop\_back}-Lock ($4.7 \to 0.0$), rear-right firing while stopped on \texttt{back\_right\_intensity}-Lock ($100 \to 0$), and backward-walking firing on \texttt{walk\_stop\_back}-Walking ($9.7 \to 0$). On the front-right event (\texttt{front\_right\_wire}), binary firing rose from $30.7\%$ to $78.5\%$ and correct-leg firing rose from $4\%$ to $100\%$. This specialization was not free: adding the hard-negative Lock and backward-walking recordings cost a small amount of accuracy on the original, easier protocol in exchange for eliminating these field false-fire modes.

\subsection{Runtime}
Finally, the detector must run inside the control loop. Table~\ref{tab:runtime} and the latency histogram in Fig.~\ref{fig:latency} report per-window inference cost for the deployment runtime under ONNX Runtime on a single CPU thread; this is a development-machine benchmark, not a measurement on the Go2's onboard computer. Mean and median latency are both about $1.0$~ms/window, with a p95 of about $1.3$~ms, comfortably under the $2.0$~ms budget implied by the $500$~Hz cadence; the exported model is about $0.5$~MB. We also verified that the deployment runtime reproduces the research pipeline numerically, with maximum absolute deviation $\leq 3.8\times10^{-7}$ (detection), $6.3\times10^{-7}$ (per-leg), and $2.1\times10^{-7}$ (intensity) over the test windows, so the offline metrics above transfer to the deployed node to within roughly $10^{-6}$.